\chapter{Device maps}
\label{ch:maps}

A device map describes the registers of a device model: their names, areas,
addresses, types, byte and word order, scaling, units and state labels, plus
the device's connection defaults. With a map, Modbench reads a device by
name, as its manual describes it, instead of as a list of numbers.

Maps are \file{.mbmap.json} files in an open format. Over 400 of them, for
meters, inverters, drives and controllers, are shared in the
\emph{mbmap-library} and can be downloaded from the Library page
(Chapter~\ref{ch:library}).

\section{Opening a map}

\begin{itemize}
  \item In step 2 of the Main page, set \field{Source} to the map, and pick it
        in \field{Map}: the maps of your library are listed there.
  \item \nav{From a file…} opens any \file{.mbmap.json} or \file{.mbmap}.
  \item Double-click a map file: it opens in Modbench.
  \item From the Library page: right-click a map and choose \nav{Use}.
\end{itemize}

Several maps can be open at once; step 2 switches between them.
The map's connection defaults (unit ID, port, baud rate, parity) fill step 1
when it is opened.

\begin{note}[Older maps]
A \file{.mbmap} file in the version 1 format opens with a warning, and
Modbench offers to convert it: the converted copy is saved next to it as
\file{.mbmap.json}, and the original is kept.
\end{note}

\section{The Device map tab}

\screenshot[8cm]{device-map}{The Device map tab during a test.}
  {The Device map tab with a meter's map: groups, names, values with units,
   some state labels, a value just changed (highlighted), the search field,
   Read all / Read none and the address menu.}

Every register of the map is a row: its name, address, value with its unit,
and its group. Values update with each cycle of the test.

\begin{description}[style=nextline, leftmargin=0pt, font=\intermedium]
  \item[Choosing what to read] Untick a register to stop reading it;
    \ui{Read all} and \ui{Read none} change them all. Fewer registers make
    each cycle shorter. Modbench joins neighbouring registers into as few
    requests as the device allows.
  \item[Search] \field{Search name or address} filters the list.
  \item[State labels] Registers that hold a state (\val{0} Stopped, \val{1}
    Running…) show its label; reserved values the map declares (such as
    \val{0xFFFF} for ``not implemented'') show their meaning instead of a
    number.
  \item[Bits and bit fields] Registers that pack several flags show each bit
    or field by its own name.
  \item[Changed values] A value that changes stays highlighted for a while
    (\nav{Settings\then Values}), so changes stand out in a long list.
  \item[Right-click] Copy a value, a name or an address, add the register to
    the chart, or write it.
\end{description}

\subsection{Registers the map is missing}

\ui{Add a register the map is missing…} adds a register on this computer
only: its address, type and name. It is read with the others and marked as
added; it is not written into the map file.

\section{Is it the right map?}

A wrong map, base address or byte order gives values that look like numbers
but mean nothing. Modbench checks in two ways:

\begin{description}[style=nextline, leftmargin=0pt, font=\intermedium]
  \item[Reference values] A map can give the value a healthy device holds in
    some registers (a model code, a fixed protocol identifier). Modbench reads
    them and shows \emph{match}, \emph{mismatch} or \emph{failed} (no answer).
    A mismatch points to a wrong map or settings.
  \item[Models and firmware] When the map says which models and firmware
    versions it fits, and has registers that identify the device, Modbench
    compares them with the connected device. \emph{Device recognised} confirms
    it; otherwise a banner says the device does not match the map. Once the
    firmware is known, registers that do not exist in that version are left
    out.
\end{description}

\screenshot[2.5cm]{map-mismatch}{A device that does not match the map.}
  {The banner on the Device map tab: the device reports another model and a
   firmware outside the map's range. \emph{Values may be wrong. Writes ask for
   confirmation first.}}

\begin{warning}
After a mismatch, writes ask for confirmation even when confirmation is off.
\end{warning}

\section{Checking a device against its map}
\label{sec:check}

When the connection and the map are set, \ui{Check device} (the tick button in
the toolbar of the Device map tab) reads every register of the map once and
gives each a result: \emph{passed}, \emph{failed} or \emph{not read}. It also
compares the model and the firmware the device reports with what the map
declares, like the checks above. The summary counts the three results, or says
\emph{No failures}.

\ui{Save report (HTML)} writes the result to one page, to keep with a
commissioning or to send. Check device needs the connection to itself, so stop
reading the map first. It only reads, and never writes to the device. The
command line does the same with \val{modbench-cli check}
(Chapter~\ref{ch:cli}).

\section{Map status}
\label{sec:map-status}

Maps from the library carry a status, shown as an icon in the library and in
the map picker:

\begin{center}
\begin{tabularx}{\linewidth}{l L}
  \toprule
  \thead{Status} & \thead{Meaning} \\
  \midrule
  Verified & Checked against a real device. \\
  Official & Written from the manufacturer's documentation, not yet
    checked on a device. \\
  Unofficial & Not from the manufacturer's documentation, for example
    reverse-engineered. \\
  \bottomrule
\end{tabularx}
\end{center}

\section{Byte order override}

\field{Byte order} in step 2 reads every value with another byte and word
order than the map says: \emph{as the map says}, or one of the four orders.
Use it to check a map whose values look wrong; it is remembered per device.

\section{Computed values}
\label{sec:computed}

\ui{Add a computed value…} under the registers adds a value worked out from
others: a name, a unit and a formula, for example power from voltage and
current, or the difference between two counters:

\begin{lstlisting}
Voltage_L1 * Current_L1 / 1000
[Energy import] - [Energy export]
round(Power / Power_max * 100, 1)
\end{lstlisting}

A register is named by its name with underscores for spaces and other signs
(\val{Voltage L1} is \val{Voltage\_L1}), or by its exact name in brackets.
Formulas take \val{+ - * / \%} and \val{**} (power), comparisons
(\val{< <= > >= == !=}, which give 1 or 0), \val{\&\& ||}, bit operators
(\val{\& | \^{} ~ << >>}), hexadecimal and binary numbers, and the functions
\val{abs}, \val{min}, \val{max}, \val{avg}, \val{round}, \val{floor},
\val{ceil}, \val{sqrt} and \val{pow}. A formula with a mistake or a name the
map does not have is refused, saying what is wrong.

A computed value is shown with \ui{calc} instead of an address and its formula
as its type. It is worked out every cycle from the values just read and can be
charted, recorded and given alarms like a register. When one of its inputs
failed, it has no value (a gap) and says which input. Like registers added by
hand, computed values are kept for the map file on this computer and never
written into it; \ui{Remove} deletes one.

\section{Exporting and recording}

\begin{description}[style=nextline, leftmargin=0pt, font=\intermedium]
  \item[\ui{Export…}] Saves the values shown, in the order shown, to a CSV
    file.
  \item[\ui{Record…}] Appends every cycle of the test to a CSV file until
    recording stops: one row per cycle, one column per register. A read
    that failed is written \val{\#failed}, and a reserved value its label
    (\val{Unimplemented}), so neither is mistaken for a value. The Trend
    player plays such a file back (Section~\ref{sec:trend-player}) with the
    same marks as the Chart tab; older recordings, which left failed reads
    blank, still open.
\end{description}

The CSV separator is set in \nav{Settings} (comma or semicolon, for
spreadsheets in languages that use a decimal comma).

\section{What is remembered per device}

For each device (map, connection and unit ID), Modbench remembers which
registers were switched off, the byte order override and Loop. They come back
when you return to that device; restoring them never writes and never starts
a test. \nav{Settings\then Connection} turns this off, and forgets what was
saved.
